\label{chap:conclusion}

\section{Summary of Findings}

This thesis asked whether a lie detector for a 70B-parameter model, built without any labelled lies,
can separate lies from honest answers when the two use the same words, while keeping the transfer of
persona directions across datasets. Our answer is one detector that combines the extended persona
contrast with the knowledge direction K. The answers to the research questions of
\cref{sec:objective} follow.

\emph{RQ1: transfer.} The selected detector, the extended contrast plus K, reaches
\V{th-k-ninek-off-L40} mean off-diagonal AUROC at layer~40 across the six datasets, against
\V{off-base-L40} for probes trained on one dataset, and its five-dataset mean of
\V{th-k-ninek-mean-L40} compares with \V{lb-mean-auroc-avg} for the benchmark's Mean Probe, with no
labels from any of them (\cref{sec:transfer,sec:selection}). Most of the transfer comes from its
persona part: the extended contrast alone reaches \V{off-ext-L40}, and the two are level on transfer
and on the five-dataset mean at this layer. The published method's trained components do not beat raw
activations at this scale.

\emph{RQ2: the same words.} On \hpkr's lies and honest answers that are the same sentence, the
selected detector reads \V{th-k-ninek-ident-L40} at layer~40, above the random band, where every
persona direction is inside it. Read through the Gender-Secret fine-tune, it detects the fine-tune's
lies at \V{th-k-gs-ft-ninekft-L40}, where the persona directions and the probes are near chance and the
benchmark's Mean Probe reads \V{th-lb-mean-auroc-gs} (\cref{sec:k-results}). Both gains come from K:
alone it reads \V{th-k-k-ident-L40} on the identical text and \V{th-k-gs-ft-kft-L40} through the
fine-tune, while the base model's read of the same transcripts gives only \V{th-k-gs-base-kft-L40}. On
Soft-Trigger its high scores cannot be credited, because the released data let the trigger itself
predict the label.

\emph{RQ3: the two parts.} The persona contrast reads how a response is framed. On \hpkr{} the
knowledge-denial contrast flags \V{th-k-deight-flag-hden-L40}\% of the honest denials at the
benchmark's threshold, every persona direction is within reach of random directions once the words are
held fixed, and all are at chance on Gender-Secret (\cref{sec:hpkr-tests}). We read this as a register
effect: a persona direction responds to how a response is framed more than to whether it is a lie
(\cref{sec:interpretation}). The knowledge direction, built from \V{th-k-known} questions the model
knows and \V{th-k-unknown} that nobody can know with the same denial forced after each, reads what the
model knows while the words stay fixed. It uses no labelled lie, shares no question with any evaluation
set, and its wording component is not what separates the identical text
(\cref{sec:k-design,sec:k-results}). Each part covers the other's weakness: K alone averages
\V{th-k-k-mean-L40} on the five datasets, below the persona contrast, and the persona contrast alone is
blind where the words are fixed.

\emph{RQ4: robustness and operating points.} On transfer the selected detector stays ahead of the
published contrast, the unified contrast and the raw probe at every one of the \V{depth-n} layers,
though at layer~\V{th-k-ninek-depth-lead-min-layer} it leads by only \V{th-k-ninek-depth-lead-min}, and
on the identical text it is inside the random band below layer~\V{th-k-ninek-depth-ident-above-from},
where K is weak. K's identical-text result survives resampling its questions, is reached by none of
1,000 splits of the same questions that remove the knowledge contrast, and replicates with new
questions written afterwards, against criteria fixed in advance, at a cosine of \V{th-k-rep-cos-L40}
with the original (\cref{sec:k-robustness}). The weakness is the tail. At 1\% of honest answers flagged
K catches \V{th-k-op-k-1-L40} of lies, and adding it to D9 lowers D9's \V{th-k-op-dnine-1-L40} to
\V{th-k-op-ninek-1-L40}; a better-placed threshold does not fix this, because what limits detection
there is the overlap between lies and the highest-scoring honest answers (\cref{sec:stats-calibration}).

\emph{Across models.} Repeated on the benchmark's three other models (\cref{sec:xmodel}), the persona
result carries over to Mistral-Small-3.1-24B and Gemma-3-27B, where the extended contrast transfers at
\V{xm-mistral-off-dnine} and \V{xm-gemma-off-dnine}, above the raw probe's \V{xm-mistral-off-raw} and
\V{xm-gemma-off-raw}, but not to Qwen2.5-72B, which played only \V{xm-qwen-roleplay}\% of its persona
rollouts fully in role. On all four models the knowledge-denial roles raise the persona contrast's mean,
and the selected detector ranks lies better than the base probe trained at the same depth. K separates
identical denials on Mistral (\V{xm-mistral-ident-k} against a band of \V{xm-mistral-ident-band}) and
narrowly on Qwen, but it does not detect the lies of Mistral's Gender-Secret fine-tune, the one other
fine-tune that could test it, and at Gemma's middle block it is weak on its own. The selected design's
lead at every layer is Llama's alone.

The selected design is the extended contrast plus the knowledge direction, read at layer~40 and combined
by a fixed rule. It is a single 8192-dimensional vector, computed once from role-play and from 120
questions, and applied with one inner product per response. Steering along its persona part, the
extended contrast, cuts disclosure of a hidden fact from \V{th-steer-base} to \V{th-steer-ext-plus}
where random directions do not, so the model uses that direction, plausibly as the register of the
deceptive roles; its knowledge part was not steered (\cref{sec:steering}).

\section{Contributions to the Field}

\begin{itemize}[leftmargin=1.2em, itemsep=2pt]
\item A label-free lie detector, the extended persona contrast plus K, that transfers better than
  supervised probes and separates lies that use the same words as honest answers.
\item A knowledge direction that reads what the model knows with the words held fixed, built without
  labelled lies.
\item A persona bank for deception by denial, with which a single persona contrast transfers at 70B
  better than supervised probes.
\item A diagnosis of persona directions: on lies about knowledge they read the words of a denial.
\item Tests on fine-tuned models that lie, where K detects Gender-Secret's lies and the persona
  directions do not.
\item A causal test of the persona part: steering along it makes the model conceal facts, unlike
  random directions.
\item A replication on the benchmark's three other models, which separates what holds across models
  (the knowledge-denial roles, the selected detector over a probe at the same depth, K on identical
  denials) from what holds for Llama (the lead at every layer, K on Gender-Secret), and shows that a
  persona direction fails on a model that does not play its roles.
\end{itemize}

\section{Recommendations for Future Work}

\begin{itemize}[leftmargin=1.2em, itemsep=2pt]
\item Other models. On the benchmark's three other models the persona contrasts fail on Qwen, whose
  role-play is the weakest, and K is weak at Gemma's middle block and does not detect Mistral's
  Gender-Secret lies (\cref{sec:xmodel}). Whether a bank that Qwen plays fully would rescue it, why K's
  knowledge contrast does not separate Gemma's lies, and whether persona and knowledge directions transfer
  between model families, are open.
\item Calibration in the setting. The knowledge direction ranks well but catches few lies at 1\% false
  positives. Cutoffs set from a setting's own honest answers help a little; directions that separate
  lies from the highest-scoring honest answers are the next step.
\item Steering along the knowledge direction. We steered only along persona directions. Adding or
  removing K while the model answers would test whether the model uses it, and whether it can make a
  model admit what it knows.
\item Other kinds of knowledge. K was built from general facts against everyday particulars. Knowledge
  of hazardous subjects, of the user, or of the model's own actions may need questions of their own.
\item Lies that are not about knowledge. K reads one ingredient of a lie. Lies that misstate intent,
  such as sandbagging or insider trading, need a second label-free signal, and a combination of K with a
  black-box auditor or a self-evaluation channel, as in Aletheia's Quest \cite{birardi2026aletheia}, may
  cover more of them.
\item Cleaner model organisms. A fine-tuned liar tests a detector only if the base model cannot already
  separate its lies from the text. Organisms whose triggers are not readable from the transcript would
  make that test possible.
\end{itemize}
